\section{Conclusion}
\label{sec:conclusion}

\methodname{} distills textured 3D morphing into two source-conditioned
flows with gated attention and a frozen prior. EPRE complements
distribution and path metrics with source-relative progression.
The model obtains the lowest EPRE under the Dense Protocol, with trade-offs in path length
and step variation. Final matched-step ablations remain pending.

\subsection{Limitations}
\label{sec:limitations}

The student inherits the coverage and biases of its teacher-generated
training data. The frozen prior supplies another constraint from TRELLIS,
not ground-truth supervision for a unique morph. Direct SLat conditioning
also remains limited by the source encoder, sparse representation, and
frozen decoders. Selecting a single source in the attention gate at an
endpoint therefore does not guarantee exact reconstruction of that asset.

Each interpolation coefficient is sampled separately. Reusing noise
across coefficients helps align the generations, but does not guarantee
continuous geometry or stable topology. Likewise, source-order
augmentation encourages reversal consistency without enforcing it
exactly. A sequence-level objective could address these constraints more
directly.

Our measurements use rendered appearance, which does not establish
geometric correctness or semantic plausibility. EPRE measures agreement
with a linear distance profile; a valid nonlinear morph can receive a
higher error. Since canonical sources are inserted as evaluation
endpoints, it also does not test endpoint reconstruction. Normalization
can amplify small distance errors when the two sources are perceptually
very similar. The present
comparison covers 50 pairs with fixed checkpoints, rather than variation
across independent training runs. Final matched-step ablations and broader
asset coverage are needed to isolate component effects and assess
generalization.
